% 注释的快捷键: Ctrl + /
% \documentclass[a4paper]{ctexart}

\documentclass[12pt,a4paper]{article}

\usepackage{xeCJK}
% \usepackage{ctex}
\usepackage{amsmath, amssymb, amsfonts, amsthm, bbm}
\usepackage{mathrsfs}
\usepackage{gensymb}
\usepackage{graphicx}
\usepackage{hyperref}
\usepackage{url}
\usepackage{cases}
\usepackage{color}
\usepackage{tikz-cd}
% \usepackage{quiver}
\usepackage{setspace}
\usepackage{spectralsequences}

\bibliographystyle{plain}

\allowdisplaybreaks
% \usepackage{times}
% \usepackage{mathptmx}
% \DeclareMathAlphabet{\mathcal}{OMS}{cmsy}{m}{n}
% \DeclareSymbolFont{largesymbols}{OMX}{cmex}{m}{n}

% \usepackage{unicode-math}
% \setmathfont{XITS Math}
% \setmathfont{XITS Math}[StylisticSet=1,range=cal]

%下面是引用前文定理的代码
% \usepackage{hyperref}
% \begin{lemma}[test]\label{test}
	%     test
	%   \end{lemma}

%   number: \ref{test}, name: \nameref{test}

%如何在等号上加文字或记号
%通过 \stackrel{?}{=}

% 如何实现大括号多编号
% \usepackage{cases}
% \begin{numcases}{}
% 	x_1&=eq1 \label{eqsystem1} \\
% 	x_2+1&=eq2 \label{eqsystem2}
% \end{numcases}

% 如何插入交换图


% 如何插入图片
% \begin{figure}[htbp]
%     \centering
%     \includegraphics[scale = ...]{name}
%     \caption{...}
% \end{figure}

% 如何插入脚注
% \footnote{}
% \footnotemark[number]
% \footnotetext[number]{text}

\newtheorem{theorem}{Theorem}[section]
\newtheorem{proposition}[theorem]{Proposition}
\newtheorem{lemma}[theorem]{Lemma}
\newtheorem{corollary}[theorem]{Corollary}
\theoremstyle{definition}
\newtheorem{definition}[theorem]{Definition}
\newtheorem{example}[theorem]{Example}
\newtheorem{problem}[theorem]{Problem}
\newtheorem{exercise}[theorem]{Exercise}
\newtheorem*{remark}{Remark}
\newtheorem*{problemstar}{Problem}
\renewcommand{\proofname}{\bf Proof}
\newcommand{\pd}[2]{\frac{\partial #1}{\partial #2}} %一次偏导
\newcommand{\npd}[3]{\frac{\partial^{#1}{#2}}{\partial{#3}^{#1}}} %n次偏导, 也可以指代混合偏导
\newcommand{\pD}[2]{\frac{{\rm D}(#1)}{{\rm D}(#2)}} %Jacobian
\newcommand{\lc}[2]{\nabla_{#1}{#2}}
\newcommand{\ft}[3]{{#1}_{#2},\dots,{#1}_{#3}}
\newcommand{\cft}[3]{{#1}_{#2},\cdots,{#1}_{#3}}
\newcommand{\id}{{\rm id}}
\newcommand{\Id}{{\rm Id}}
\newcommand{\supp}{{\rm supp}}
\newcommand{\mi}{\sqrt{-1}}
\newcommand{\me}{{\rm e}}
\newcommand{\md}{{\rm d}}
\newcommand{\mZ}{\mathbb{Z}}
\newcommand{\mQ}{\mathbb{Q}}
\newcommand{\mR}{\mathbb{R}}
\newcommand{\mC}{\mathbb{C}}
\newcommand{\GL}{{\rm GL}}
\newcommand{\SL}{{\rm SL}}
\newcommand{\Tr}{{\rm Tr}}
\newcommand{\str}{{\rm Tr}_s}
\newcommand{\Aut}{{\rm Aut}}
\newcommand{\Inn}{{\rm Inn}}
\newcommand{\Out}{{\rm Out}}
\newcommand{\End}{{\rm End}}
\newcommand{\Lie}{{\rm Lie}}
\newcommand{\mCP}{\mathbb{C}\mathrm{P}}
\newcommand{\mRP}{\mathbb{R}\mathrm{P}}
\newcommand{\mHP}{\mathbb{H}\mathrm{P}}
\newcommand{\vep}{\varepsilon}
\newcommand{\ra}{\rightarrow}
\newcommand{\gs}{\geqslant}
\newcommand{\ls}{\leqslant}
\newcommand{\mbb}{\mathbb}
\newcommand{\mc}{\mathcal}
\newcommand{\mfr}{\mathfrak}
\newcommand{\mrm}{\mathrm}
\newcommand{\Hom}{\mathrm{Hom}}
\newcommand{\Ext}{\mathrm{Ext}}
\newcommand{\Tor}{\mathrm{Tor}}
\newcommand{\Ker}{\mathrm{Ker}}
\newcommand{\Gr}{\mathrm{Gr}}
\newcommand{\p}{\partial}
\newcommand{\bp}{\bar{\partial}}

% \graphicspath{{Figures/}}

\begin{document}

\title{Characteristic Classes\\
    \large Book Notes of \textit{Characteristic Classes}}
\author{Jianbo Zhou}
\date{\today}
\maketitle
\tableofcontents
\newpage

% \section{Obstruction Theory \& Characteristic Classes}
\label{sec:ObstructionTheory&CharacteristicClasses}
\subsection{A Brief Introduction to Obstruction Theory}
In homotopy theory, one often encounters the following problem: 
\begin{center}
\fbox{
    \parbox{0.9\linewidth}{
        Given a pair of spaces $(X,A)$ and a map $f:A\to Y$. Can $f$ be extended 
        to a map $F:X\to Y$?
    }
}
\end{center}
Obstruction theory provides a systematic way to study 
this problem. 

$\mathbf{Toy\ Model:}$ A map $f:S^n\to Y$ can be extended to a map 
$F:D^{n+1}\to Y$ if and only if $f$ is null-homotopic.

Now we consider the general case. Let $K$ be a CW complex, $Y$ be a simple 
space, i.e., $\pi_1(Y)$ acts trivially on $\pi_n(Y)$ for all $n$. Under this 
condition, we can unambiguously use the notation $\pi_n(Y) = [S^n, Y]$.  
Suppose we have a map $f:K^{(n)}\to Y$, where $K^{(n)}$ is the 
$n$-skeleton of $K$. We want to extend $f$ to a map $F:K^{(n+1)}\to Y$. 
To do this, we only need to extend $f$ to each $(n+1)$-cell $e_\alpha^{n+1}$. 

Suppose the attaching map of $e_\alpha^n$ is 
$\varphi_\alpha:S^{n}\to K^{(n)}.$ Then $f\circ \varphi_\alpha$ is a map
from $S^{n}$ to $Y$. By the toy model, $f\circ \varphi_\alpha$ can be
extended to $D^{n+1}$ if and only if $[f\circ \varphi_\alpha] = 0$ in
$\pi_{n}(Y)$. We in fact get a cochain 
$c^{n+1}(f)\in C^{n+1}_{cell}(K;\pi_{n}(Y))$: 
\begin{align*}
    c^{n+1}(f):C_{n+1}^{cell}(K)\to \pi_{n}(Y) \\
    e_\alpha^{n+1}\mapsto [f\circ \varphi_\alpha].
\end{align*}
We call $c^{n+1}(f)$ the \textbf{${(n+1)}$-th obstruction cochain} of $f$. 
There are some basic facts about $c^{n+1}(f)$: 
\begin{itemize}
    \item $f$ can be extended to $K^{(n+1)}$ if and only if $c^{n+1}(f) = 0$; 
    \item If $f,f':K^{(n)}\to Y$ are homotopic, then 
    $c^{n+1}(f) = c^{n+1}(f')$; 
    \item (Natruality) Suppose $L$ is another CW complex and $\phi:L\to K$ 
    is a cellular map, then $c^{n+1}(f\circ \phi) = \phi^{\sharp}c^{n+1}(f)$.
\end{itemize}

At the following, we give some important properties of $c^{n+1}(f)$ without 
proof. 
\begin{theorem}
    $c^{n+1}(f)$ is a cocycle, i.e., $\delta c^{n+1}(f) = 0$. Thus it 
    defines a cohomology class $[c^{n+1}(f)]\in H^{n+1}(K;\pi_{n}(Y))$. 
    We call $[c^{n+1}(f)]$ the \textbf{$(n+1)$-th obstruction class} of $f$.
\end{theorem}

A similar concept is defined for homotopies. Suppose we have two maps
$f_0,f_1:K^{(n)}\to Y$ and a homotopy $h:K^{(n-1)}\times I\to Y$ between
$f_0|_{K^{(n-1)}}$ and $f_1|_{K^{(n-1)}}$. We want to extend $h$ to a 
homotopy $H:K^{(n)}\times I\to Y$ between $f_0$ and $f_1$. To do this, 
we only need to extend $h$ to each $n$-cell $e_\alpha^n$. Let 
$\varphi_\alpha:S^{n-1}\to K^{(n-1)}$ be the attaching map of $e_\alpha^n$, 
$\psi_\alpha:D^n\to K^{(n)}$ be the characteristic map of $e_\alpha^n$. Then 
we can form a map
\begin{align*}
    h_\alpha:(D^n\times \{0,&1\})\cup(\p D^n\times I)\to Y \\ 
    (x,0)&\mapsto f_0\circ \psi_\alpha(x) \\
    (x,1)&\mapsto f_1\circ \psi_\alpha(x) \\
    (x,t)&\mapsto h\circ (\varphi_\alpha(x),t).
\end{align*}
Since $(D^n\times \{0,1\})\cup(\p D^n\times I)\cong S^n$, $h_\alpha$ defines
a map from $S^n$ to $Y$. It can extend to $D^n\times I\cong D^{n+1}$ if and 
only if $[h_\alpha] = 0$ in $\pi_n(Y)$. We in fact get a cochain
$c^n(f_0,f_1,h)\in C^n_{cell}(K;\pi_n(Y))$:
\begin{align*}
    c^n(f_0,f_1,h):C_n^{cell}(K) & \to \pi_n(Y) \\
    e_\alpha^n & \mapsto [h_\alpha]. 
\end{align*} 
We call $c^n(f_0,f_1,h)$ the \textbf{$n$-th deformation cochain} of 
$(f_0,f_1,h)$. If $f_0|_{K^{(n-1)}} = f_1|_{K^{(n-1)}}$ and $h$ is the 
trivial homotopy, we simply write $c^n(f_0,f_1)$ for $c^n(f_0,f_1,h)$ and 
call it the \textbf{$n$-th difference cochain} of $(f_0,f_1)$. 
There are some basic facts about $c^n(f_0,f_1,h)$:
\begin{itemize}
    \item $h$ can be extended to a homotopy $H:K^{(n)}\times I\to Y$ between 
    $f_0$ and $f_1$ if and only if $c^n(f_0,f_1,h) = 0$; 
    % \item If $h,h':K^{(n-1)}\times I\to Y$ are homotopic rel. boundary, then 
    % $c^n(f_0,f_1,h) = c^n(f_0,f_1,h')$; 
    \item (Natruality) Suppose $L$ is another CW complex and $\phi:L\to K$ 
    is a cellular map, then 
    $c^n(f_0\circ \phi,f_1\circ \phi,h\circ(\phi\times \id_I)) = 
    \phi^{\sharp}c^n(f_0,f_1,h)$.
    \item $c^n(f_0,f_2,h*h') = c^n(f_0,f_1,h) + c^n(f_1,f_2,h')$, where
    $h*h'$ is the concatenation of $h$ and $h'$. Especially, if 
    $f_0|_{K^{(n-1)}} = f_1|_{K^{(n-1)}} = f_2|_{K^{(n-1)}}$ and $h,h'$ are
    trivial homotopies, then $c^n(f_0,f_2) = c^n(f_0,f_1) + c^n(f_1,f_2)$.
\end{itemize}

In general, $c^n(f_0,f_1,h)$ is not a cocycle. But we have the following
formula. 
\begin{theorem}
    $\delta c^n(f_0,f_1,h) = c^{n+1}(f_1) - c^{n+1}(f_0)$.
\end{theorem} 

There is a useful lemma about difference cochains.
\begin{lemma}
    For any $f:K^{(n)}\to Y$ and any cochain $d\in C^n_{cell}(K;\pi_n(Y))$, 
    there exists a map $f':K^{(n)}\to Y$ such that 
    $f'|_{K^{(n-1)}} = f|_{K^{(n-1)}}$ and $d^n(f,f') = d$.
\end{lemma}

With the above preparations, we have the following theorem:
\begin{theorem}[Eilenberg]
    Given a map $f:K^{(n)}\to Y$, $[c^n(f)] = 0$ in 
    $H^{n+1}(K;\pi_n(Y))$ if and only if $f$ can extend to $K^{(n+1)}$ 
    after having been changed on $K^{(n)}\big\backslash K^{(n-1)}$.
\end{theorem}
\begin{proof}
    $(\Leftarrow):$ If $\exists\,f':K^{(n+1)}\to Y$ such that 
    $f'|_{K^{(n-1)}} = f|_{K^{(n-1)}}$, then 
    \begin{equation*}
        c^{n+1}(f) = c^{n+1}(f) - c^{n+1}(f') = \delta d^n(f',f)
    \end{equation*}
    Thus $c^{n+1}(f)$ is a coboundary. 

    $(\Rightarrow):$ If  $[c^{n+1}(f)]=0$, then $\exists\,d\in C^n(K;\pi_n(Y))$ 
    such that $c^{n+1}(f) = \delta d$. By lemma, there exists 
    $f':K^{(n)}\to Y$ such that $f|_{K^{(n-1)}} = f'|_{K^{(n-1)}}$ and 
    $d^n(f,f') = -d$. So 
    \begin{equation*}
        c^{n+1}(f') = c^{n+1}(f) + \delta d^n(f,f') = \delta d - \delta d = 0.
    \end{equation*}
    Thus $f'$ can be extended to $K^{(n+1)}$.
\end{proof}

\begin{example}
    If the target space $Y$ is $n$-connected, i.e., 
    $\pi_k(Y)=0,\; k=0,1,\dots,n$, then every map $f:K^{(0)}\to Y$ can be 
    extended to $K^{(n+1)}$. Besides, any two maps $f,g:K^{(n)}\to Y$ are 
    homotopic. Indeed, $f|_{K^{(0)}}\simeq g|_{K^{(0)}}$ for the reason that 
    $K$ is path-connected. We denote the homotopy by $h_0$. Then $h_0$ can 
    be extened to $h_n:K^{(n)}\times I\to Y$ since all the deformation cochain 
    vanish.
\end{example}

\begin{corollary}
    If $Y$ is $n$-connected and $\dim K\leq n$, then all the maps from $K$ 
    to $Y$ are null-homotopic.
\end{corollary}

\begin{theorem}[Hopf-Whitehead]
    For any CW complex $K$ with $\dim K\leq n$ and $(n-1)$-connected space $Y$, 
    there is a one-to-one correspondence:
    \begin{equation*}
        [K,Y]\longleftrightarrow H^n(K;\pi_n(Y)).
    \end{equation*}
\end{theorem}
If we weaken the dimensional condition and instead strengthen the target space 
to be an Eilenberg-Maclane space $K(\pi,n)$. Then we will get another theorem: 
\begin{theorem}
    For any CW complex $X$, abelian group $\pi$ and integer $n$, there is a 
    one-to-one correspondence: 
    \begin{equation*}
        [X,K(\pi,n)]\longleftrightarrow H^n(X;\pi).
    \end{equation*}
\end{theorem}
\begin{remark}
    In fact, there is a special class $\alpha$ in 
    \begin{equation*}
        H^n(K(\pi,n);\pi)\cong \Hom(H_n(K(\pi,n)),\pi)\cong \Hom(\pi,\pi)
    \end{equation*}
    which corresponds to the identity map in $\Hom(\pi,\pi)$. The correspondence 
    can be explicitly written by 
    \begin{align*}
        [X,K(\pi,n)]&\to H^n(X;\pi) \\
        f&\mapsto f^*\alpha.
    \end{align*}
\end{remark}

\subsection{Characteristic Classes via Obsruction Theory}
\subsubsection{Cohomology with Local Coefficients}
To discrib characteristic classes as obstruction classes, we need cohomology 
with local coefficients. 
% \begin{center}
%     \fbox{
%         \parbox{0.95\linewidth}{
%             Suppose $X$ is a path-connected space. A \textbf{local coefficient system} 
%             on $X$ is a functor
%             \begin{equation*}
%                 \mathcal{G}:\pi(X)\to \text{(Abelian groups)},
%             \end{equation*}
%             where $\pi(X)$ is the fundamental groupoid of $X$. If we fix a base point
%             $x_0\in X$, then $\mathcal{G}$ is determined by the group 
%             $G = \mathcal{G}(x_0)$ and a representation $\rho:\pi_1(X,x_0)\to \Aut(G)$. 
%             We call $G$ the \textbf{typical group} of $\mathcal{G}$.
%         }
%     }
% \end{center}

\noindent\textbf{Local Coefficient System:}

Suppose $X$ is a path-connected space. For any $x\in X$, let $G_x$ be an 
abelian group attaching to $x$. For any path $\gamma_{xy}$ from $x$ to $y$, 
we have an isomorphism $(\gamma_{xy})_*:G_x\to G_y$. If $(\gamma_{xy})_*$ 
only depends on the homotopy class of $\gamma$ rel. endpoints and satisfies 
the following conditions: 
\begin{equation*}
    (\text{constant path})_* = \Id, \quad
    (\gamma_{xy}*\gamma_{yz})_* = (\gamma_{yz})_*\circ(\gamma_{xy})_*.
\end{equation*}
Then we call $\mathcal{G} = \{G_x\}_{x\in X}$ a 
\textbf{local coefficient system} on $X$. If we fix a base point $x_0\in X$, 
then $\mathcal{G}$ is determined by the group $G = G_{x_0}$ and an action 
$\rho:\pi_1(X,x_0)\to \Aut(G)$. We call $G$ the \textbf{typical group} of 
$\mathcal{G}$. 

\begin{example}
    If $\mathcal{G}$ is a local coefficient system on $X$ with typical group 
    $G$ and trivial action. Then it is the ordinary constant coefficient we 
    are familiar with.
\end{example}

\begin{example}
    If $G_x = \mZ_2$ for all $x\in X$, then since $\Aut(\mZ_2)=\{\Id\}$, 
    the local coefficient system with typical group $\mZ_2$ must be trivial. 
\end{example}

\begin{example}
    If $G_x = \mZ$ for all $x\in X$ and $\pi_1(X,x_0)$ acts on $\mZ$ by
    multiplication of $\pm 1$ according to whether the orientation is
    preserved or reversed along the loop. Then $\mathcal{G}$ is called the 
    \textbf{orientation system} on $X$ and denoted by $\mathrm{Or}$. 
\end{example}

\begin{example}
    Let $p:\tilde{X}\to X$ be the universal covering of $X$. Then 
    $\mathcal{G} = \{H_n(p^{-1}(x))\}_{x\in X}$ is a local coefficient system 
    on $X$ with typical group $H_n(\tilde{X})$ and action induced by the 
    deck transformation group $\pi_1(X,x_0)$. 
\end{example}

\begin{example}
    Let $E\to X$ be a vector bundle with fiber $F$. Then 
    $\mathcal{G} = \{H_n(p^{-1}(x))\}_{x\in X}$ is a local coefficient system 
    on $X$ with typical group $H_n(F)$ and action induced by the structure 
    group of $E$.
\end{example}

\begin{example}
    $\mathcal{G} = \{\pi_n(X,x)\}_{x\in X}$ is a local coefficient system 
    on $X$ with typical group $\pi_n(X,x_0)$ and action induced by the 
    fundamental group $\pi_1(X,x_0)$.
\end{example}

\noindent\textbf{Cohomology with Local Coefficients:}

Let $X$ be a CW complex and $\mathcal{G} = \{G_x\}_{x\in X}$ be a local
coefficient system on $X$. We define the cellular cochain groups with local
coefficients by
\begin{equation*}
    C^n_{cell}(X;\mathcal{G}) = \left\{f:\big\{\text{$n$-cells of $X$}\big\}
    \to \bigsqcup_{x\in X}G_x \,\Big|\, f(e_\alpha^n)\in 
    G_{x_\alpha}\right\}
\end{equation*}
where $x_\alpha$ is a chosen point in the $n$-cell $e_\alpha^n$. The 
coboundary operator 
$\delta:C^n_{cell}(X;\mathcal{G})\to C^{n+1}_{cell}(X;\mathcal{G})$
is defined by 
\begin{equation*}
    (\delta f)(e_\beta^{n+1}) = \sum_\alpha d_{\beta\alpha}
    (\gamma_{\alpha\beta})_*f(e_\alpha^n),
\end{equation*}
where $d_{\beta\alpha}$ is the degree of the map 
$\p e_\beta^{n+1}\to X^{(n)}/(X^{(n)}-e^n_\alpha)\cong S^n$ and 
$\gamma_{\alpha\beta}$ is a path from $x_\alpha$ to $x_\beta$. It can be
verified that $\delta^2=0$. Thus we can define the cohomology groups with 
local coefficients by 
\begin{equation*}
    H^n(X;\mathcal{G}) = \frac{\ker(\delta:C^n_{cell}(X;\mathcal{G})\to 
    C^{n+1}_{cell}(X;\mathcal{G}))}{\mathrm{im}(\delta:C^{n-1}_{cell}(X;
    \mathcal{G})\to C^n_{cell}(X;\mathcal{G}))}.
\end{equation*}

\subsubsection*{The Homotopy Groups of $V_k(\mR^n)$}
\textbf{Fact:} The Stiefel manifold $V_k(\mR^n)$ is $(n-k-1)$-connected and 
the first nontrivial homotopy group is 
\begin{equation*}
    \pi_{n-k}(V_k(\mR^n)) = 
    \begin{cases}
        \mZ,& n-k \text{ is even or } k=1; \\
        \mZ_2,& n-k \text{ is odd and } k>1.
    \end{cases}
\end{equation*}
\begin{remark}
    It is evident that $V_k(\mR^n)$ is a simple space.
\end{remark}

\subsubsection*{Characteristic Classes}
Now let's back to vector bundles. Let $\xi$ be an $\mR^n$-bundle over a CW 
complex $B$. We want to investigate how nontrivial $\xi$ is. To do this, we 
try to find $k$-linearly independent sections of $\xi$. 

% First suppose there 
% is a Euclidean metric $g$ on $\xi$. Then we can talk about orthonormal 
% sections. 
% \begin{remark}
%     Any real vector bundle over a paracompact space admits a 
%     Euclidean metric. Each two Euclidean metrics $g_0, g_1$ are 
%     homotopic through $g_t = (1-t)g_0 + tg_1$. 
% \end{remark}

Let $V_k(\xi)$ be the Stiefel manifold bundle associated to $\xi$. 
It's a fiber bundle over $B$ with fiber $F=V_k(\mR^n)$. A section $s$ of 
$V_k(\xi)$ is a $k$-frame of $\xi$. Locally on a trivializing
neighborhood $U_\alpha$, $s$ can be identified as a map from $U_\alpha$ to
$V_k(\mR^n)$. Thus we can use obsruction theory to study the existence of
sections of $V_k(\xi)$. 

The obstructions $c^{j+1}(s)$ lie in $H^{j+1}(B;\mathcal{G})$, where 
$\mathcal{G}$ is a local coefficient system on $B$ with typical group 
$\pi_j(V_k(\mR^n))$. This is because although the ristriction of $V_k(\xi)$ 
on every $j$-cell $e^j_\alpha$ is trivial, the transition function may not 
be identity, from which induces a nontrivial isomorphism of homotopy groups. 
Moreover, any path $\gamma_{xy}$ from $x$ to $y$ can lift to a homotopy 
from $F_x$ to $F_y$, thus induces an isomorphism 
\begin{equation*}
    (\gamma_{xy})_*:\pi_j\big(F_x,x\big)
    \overset{\cong}{\to}\pi_j\big(F_y,y\big).
\end{equation*}
The theorems about obstruction classes can all generalize to cohomology with 
local coefficients without difficulty. 

Now since $V_k(\mR^n)$ is $(n-k-1)$-connected, all the section $s$ can be 
extended to $B^{(n-k)}$, all the sections on $B^{(n-k)}$ are homotopic with 
each other. The first obstruction comes when extend $s$ to 
$(n-k+1)$-skeleton, which we denote by 
\begin{equation*}
    \mathfrak{o}_{n-k+1}:=c^{n-k+1}(s)\in 
    H^{n-k+1}\Big(B;\big\{\pi_{n-k}\big(V_k(\mR^n)\big)\big\}\Big)
\end{equation*}
Since all sections on $B^{(n-k)}$ are homotopic with each other, this 
obstruction class does not depend on the choice of $s$, i.e, 
$\mathfrak{o}_{n-k+1}$ it is well-defined. 

Now we construct a series of characteristic classes from obsruction classes. 
\begin{itemize}
    \item If $n-k=2j-1$ is odd and $1<k<n$, then 
    $\pi_{2j-1}(V_{n-2j+1}(\mR^n)) = \mZ_2$. In this case, 
    \begin{align*}
        H^{2j}\Big(B;\big\{\pi_{2j-1}\big(V_{n-2j+1}(\mR^n)\big)\big\}\Big)
        &\cong H^{2j}(B;\mZ_2) \\ 
        \mathfrak{o}_{2j} &\mapsto w_{2j}(\xi).
    \end{align*}
    \item If $n-k=2j$ is even and $1<k<n$, then 
    $\pi_{2j}(V_{n-2j}(\mR^n)) = \mZ$. After applying the mod 2 reduction, 
    we have 
    \begin{align*}
        H^{2j+1}\Big(B;\big\{\pi_{2j}\big(V_{n-2j}(\mR^n)\big)\big\}\Big)
        &\xrightarrow{\text{mod } 2} H^{2j+1}(B;\mZ_2) \\ 
        \mathfrak{o}_{2j+1} &\mapsto w_{2j+1}(\xi).
    \end{align*}
    \item If $k=n$, then $V_n(\mR^n) = O(n)$, but $\pi_0(O(n)) = \mZ_2$ is 
    not a group. We turn to use $\tilde{H}^0(O(n);\mZ) = \mZ$ instead. Then 
    $\tilde{H}^0(O(n);\mZ)$ is a local coefficient system on $B$ with typical 
    group $\mZ$ and action induced by the determinant representation 
    $\det:\pi_1(B)\to O(n)\to \{\pm 1\}$. After applying the mod 2 reduction, 
    we have
    \begin{align*}
        H^1\Big(B;\big\{\tilde{H}^0(O(n))\big\}\Big)
        &\xrightarrow{\text{mod } 2} H^1(B;\mZ_2) \\ 
        \mathfrak{o}_1 &\mapsto w_1(\xi).
    \end{align*}
    \item If $k=1$, then $\pi_{n-1}(V_1(\mR^n)) = \pi_{n-1}(S^{n-1}) = \mZ$.
    After applying the mod 2 reduction, we have
    \begin{align*}
        H^n\Big(B;\big\{\pi_{n-1}(V_1(\mR^n))\big\}\Big)
        &\xrightarrow{\text{mod } 2} H^n(B;\mZ_2) \\ 
        \mathfrak{o}_n &\mapsto w_n(\xi).
    \end{align*}
\end{itemize}

\begin{center}
\begin{table}[htbp]
\begin{tabular}{|c|cc|}
\hline
Cases              & \multicolumn{1}{c|}{Obstruction classes}       & Characteristic classes       \\ \hline
$n-k=2j-1,\;1<k<n$ & \multicolumn{2}{c|}{$\mathfrak{o}_{2j}=w_{2j}$}                               \\ \hline
$n-k=2j,\;1<k<n$   & \multicolumn{2}{c|}{$\mathfrak{o}_{2j+1}\xrightarrow{\text{mod } 2}w_{2j+1}$} \\ \hline
$k=n$              & \multicolumn{2}{c|}{$\mathfrak{o}_{1}\xrightarrow{\text{mod } 2}w_{1}$}       \\ \hline
$k=1$              & \multicolumn{2}{c|}{$\mathfrak{o}_{n}\xrightarrow{\text{mod } 2}w_{n}$}       \\ \hline
\end{tabular}
\end{table}
\end{center}

\subsubsection*{Properties of our construction}
\noindent\textbf{Stiefel-Whitney Classes}
\begin{theorem}
    $w_i(\xi)\in H^i(B;\mZ_2)$ are the Stiefel-Whitney classes of $\xi$.
\end{theorem}
\begin{proof}
    omitted.
\end{proof}

\noindent\textbf{$\mathfrak{o}_n$ and the Euler Class}
If $\xi$ is orientable, then the structure group of $\xi$ can be 
reduced to $SO(n)$. Thus the action of $\pi_1(B)$ on 
$\pi_{n-k}(V_k(\mR^n))$ is trivial for $1\leq k<n$. So the local 
coefficient system is trivial. In this case, 
\begin{equation*}
    H^n\Big(B;\big\{\pi_{n-1}(V_1(\mR^n))\big\}\Big) 
    \cong H^n(B;\mZ).
\end{equation*}
\begin{theorem}
    If $\xi$ is orientable, then $\mathfrak{o}_n\in H^n(B;\mZ)$ 
    is the Euler class $e(\xi)$.
\end{theorem}
\begin{proof}
    omitted.
\end{proof}

\noindent\textbf{$w_1$ and Orientability}
The structure group of $\xi$ can be reduced to $SO(n)$ if and only if
the transition functions all have determinant 1, i.e., the action of 
$\pi_1(B)$ on $\tilde{H}^0(O(n);\mZ)$ is trivial. This is equivalent to
the triviality of $\mathfrak{o}_1$. Thus
\begin{proposition}
    $\mathfrak{o}_1 = 0$ if and only if $\xi$ is orientable.
\end{proposition}
\begin{theorem}
    $\mathfrak{o}_1$ has order 2, i.e., $2\mathfrak{o}_1 = 0$. Thus 
    \begin{equation*}
        w_1(\xi) = 0 \text{ if and only if } \mathfrak{o}_1 = 0.
    \end{equation*}
\end{theorem}
\begin{proof}
    The connected conponents of the space $V_n(\mR^n)$ are identified with 
    orientation of the fiber, and hence the elements of 
    $\tilde{H}^0(V_n(\mR^n);\mZ)\cong\mZ$ has the form 
    $a\mathrm{Or}_1-a\mathrm{Or}_2$; to each orientation we assign its 
    Coefficient.
    
    ...
\end{proof}

\noindent\textbf{Is $w_i$ the true obsruction?}
The fibration 
\begin{equation*}
    S^{2j-1}\to V_{n-2j+1}(\mR^n)\to V_{n-2j}(\mR^n).
\end{equation*}
induces an exact sequence of homotopy groups
\begin{equation*}
    \pi_{2j}(V_{n-2j}(\mR^n)) \xrightarrow{\partial} 
    \pi_{2j-1}(S^{2j-1})\to \pi_{2j-1} (V_{n-2j+1}(\mR^n)) \to 0.
\end{equation*}
It is actually 
\begin{equation*}
    \mZ \xrightarrow{\p} \mZ \xrightarrow{i_*} \mZ_2 \to 0.
\end{equation*}
Therefore $\mathrm{Im}\,\p=\Ker i_* = 2\mZ$ and $\Ker\p = 0$. So $\p$ is
multiplication by 2. So the exact sequence is a short exact sequence
\begin{equation*}
    0\to \mZ \xrightarrow{\times 2} \mZ \xrightarrow{i_*} \mZ_2 \to 0.
\end{equation*}
This short exact sequence induces a long exact sequence of cohomology groups
\begin{equation*}
    \cdots \to H^{2j}(B;\mZ) 
    \xrightarrow{i_*} H^{2j}(B;\mZ_2) \xrightarrow{\beta} H^{2j+1}(B;\mZ) 
    \xrightarrow{\times 2} H^{2j+1}(B;\mZ) \to \cdots
\end{equation*}
where $\beta$ is the Bockstein homomorphism.
\begin{theorem}
    $\mathfrak{o}_{2j+1} = \beta(\mathfrak{o}_{2j})$.
\end{theorem}
\begin{proof}
    omitted.
\end{proof}
\begin{corollary}
    $\mathfrak{o}_{2j+1}$ has order 2, i.e., 
    $2\mathfrak{o}_{2j+1} = 0$. Thus 
    \begin{equation*}
        w_{2j+1}(\xi) = 0 \text{ if and only if } 
        \mathfrak{o}_{2j+1} = 0.
    \end{equation*}
\end{corollary}
In summary, almost all the Stiefel-Whitney classes $w_i$ are the true 
obstructions, except when $i=n$ is even. If in addition $\xi$ is oriented, 
then $\mathfrak{o}_n$ is Euler class $e(\xi)$.
% \section{Chern Numbers \& Pontrjagin Numbers}
\label{sec:ChernNumbers&PontrjaginNumbers}

\subsection{Chern Numbers}
\begin{definition}
    Let $K^n$ be a compact complex manifold, 
    $\dim_\mathbb{C}K=n$. 
    For each partition $I=i_1,\dots,i_r$ of $n$, 
    we define the $I$-th Chern number to be the integer 
    \begin{equation*}
        c_I[K^n]=c_{i_1}\cdots c_{i_r}[K^n]
        :=\langle c_{i_1}(\tau)\cdots c_{i_r}(\tau),\mu_{K^n}\rangle,
    \end{equation*}
    where $\mu_{K^n}$ is the fundamental homology class w.r.t. 
    the orientation determined by the complex structure.

    If $I$ is a partition of other number, we define $c_I[K^n]$ to be $0$. 
\end{definition}

\begin{example}
    For $K^n=\mathbb{C}P^n$, we have
    \begin{align*}
        c_I[\mathbb{C}P^n] 
        &= \langle c_{i_1}(\tau)\cdots c_{i_r}(\tau),
        \mu_{\mathbb{C}P^n}\rangle \\
        &=\left\langle\binom{n+1}{i_1}\cdots\binom{n+1}{i_r}a^n,
        \mu_{\mathbb{C}P^n}\right\rangle \\
        &=\binom{n+1}{i_1}\cdots\binom{n+1}{i_r}.
    \end{align*}
\end{example}

\subsection{Pontrjagin Numbers}
\begin{definition}
    Let $M^{4n}$ be a compact oriented smooth manifold of dimension $4n$. 
    For each partition $I=i_1,\dots,i_r$ of $n$, 
    we define the $I$-th Pontrjagin number to be the integer 
    \begin{equation*}
        p_I[M^{4n}]=p_{i_1}\cdots p_{i_r}[M^{4n}]
        :=\langle p_{i_1}(\tau)\cdots p_{i_r}(\tau),\mu_{M^{4n}}\rangle,
    \end{equation*}
    where $\mu_{M^{4n}}$ is the fundamental homology class w.r.t. 
    the given orientation.

    If $I$ is a partition of other number, 
    we define $p_I[M^{4n}]$ to be $0$.
\end{definition}

\begin{example}
    For $M^{4n}=\mathbb{C}P^{2n}$, we have
    \begin{align*}
        p_I[\mathbb{C}P^{2n}] 
        &= \langle p_{i_1}(\tau)\cdots p_{i_r}(\tau),
        \mu_{\mathbb{C}P^{2n}}\rangle \\
        &=\left\langle\binom{2n+1}{i_1}\cdots\binom{2n+1}{i_r}a^{2n},
        \mu_{\mathbb{C}P^{2n}}\right\rangle \\
        &=\binom{2n+1}{i_1}\cdots\binom{2n+1}{i_r}.
    \end{align*}
\end{example}

\noindent\textbf{Oriented reversing map.} 
If we reverse the orientation of $M^{4n}$, since the definition of
Pontrjagin classes do not depend on the orientation, 
we have the following
\begin{gather*}
    p_k(M)\mapsto p_k(M), \\
    \mu_M\mapsto -\mu_M, \\
    p_I[M]\mapsto -p_I[M].
\end{gather*}
Thus the sign of Pontrjagin numbers change when 
reversing the orientation of $M$. 
If $M$ has a nonzero Pontrjagin number, 
$M$ cannot possess any orientation reversing diffeomorphism. 

\noindent\textbf{Manifold as a boundary.} 
If $M^{4n}$ is a boundary of some smooth, 
compact, oriented $(4n+1)$-dimensional manifold-with-boundary. 
Then all Pontrjagin numbers of $M$ must be zero. 
Thus if some $p_I[M]\neq0$, 
then $M$ cannot be the boundary of such manifold.

\subsection{Symmetric Functions}
We know that all symmetric polynomials in 
$\mathbb{Z}[t_1,\dots,t_n]$ is itself a polynomial ring: 
\begin{equation*}
    \mathcal{S} = \mathbb{Z}[\sigma_1,\dots,\sigma_n]\subset
    \mathbb{Z}[t_1,\dots,t_n],
\end{equation*}
where $\sigma_i$ is the $i$-th elementary symmetric polynomial.
\begin{equation*}
    1+\sigma_1+\cdots+\sigma_n = (1+t_1)\cdots(1+t_n).
\end{equation*}
If we focus on the symmetric polynomials of degree $k$, 
denoted by $\mathcal{S}^k$, clearly 
it is a free abelian group with basis
\begin{equation*}
    \sigma_{i_1}\cdots\sigma_{i_r},\quad
    I=i_1,\dots,i_r\text{ is a partition of } k.
\end{equation*}
For some reason, it is more convenient to use another basis of 
$\mathcal{S}^k$, which is described as follows.

Define two monomials in $\mathbb{Z}[t_1,\dots,t_n]$ to be 
equivalent if they differ by a permutation of the variables.
We define 
\begin{equation*}
    \sum t_1^{a_1}\cdots t_r^{a_r}
\end{equation*}
to be the summation of all distinct monomials equivalent to
$t_1^{a_1}\cdots t_n^{a_n}$. Clearly, it is a symmetric polynomial.

\begin{lemma}
    \begin{equation*}
        \sum t_1^{a_1}\cdots t_r^{a_r}
    \end{equation*}
    where $a_1\cdots a_r$ ranges over all partitions 
    of $k$ with length $r\leq n$,
    form another basis of $\mathcal{S}^k$.
\end{lemma}

Now for each partition $I=i_1,\dots,i_r$ of $k$, 
we define a polynomial $s_I$ in $k$ variables as follows:
choose $n\geq k$, define 
\begin{equation*}
    s_I(\sigma_1,\dots,\sigma_k)
    :=\sum t_1^{i_1}\cdots t_r^{i_r},
\end{equation*}
both sides are symmetric polynomials in $n$ variables.

Notes that $s_I$ does not depend on the choice of $n$, 
for if we choose another $n'>n$, 
then we can set $t_{n+1}=\cdots=t_{n'}=0$ to get back to the previous case.

Clearly, all $s_I$ form a basis of $\mathcal{S}^k$ as $I$ ranges over
all partitions of $k$.

Now let's look at the cohomology ring 
$H^*(\mathrm{Gr}_n(\mathbb{C}^\infty);\mathbb{Z})$.
the splitting map 
\begin{equation*}
    \phi:\mathbb{C}P^\infty\times\cdots\times\mathbb{C}P^\infty
    \ra \mathrm{Gr}_n(\mathbb{C}^\infty)
\end{equation*}
induces an injective homomorphism of cohomology rings
\begin{equation*}
    \phi^*:H^*(\mathrm{Gr}_n(\mathbb{C}^\infty);\mathbb{Z})
    \ra H^*(\mathbb{C}P^\infty\times\cdots\times\mathbb{C}P^\infty;
    \mathbb{Z})
    \cong \mathbb{Z}[t_1,\dots,t_n],
\end{equation*}
where $t_i$ is the generator of the $i$-th factor.
$\phi^*$ maps 
$H^*(\mathrm{Gr}_n(\mathbb{C}^\infty);\mathbb{Z})=\mathbb{Z}[c_1,\dots,c_n]$ 
isomorphically onto the symmetric polynomial subring 
$\mathcal{S}\subset\mathbb{Z}[t_1,\dots,t_n]$, 
with $c_i$ mapped to $\sigma_i$.

Since $\phi^*s_I(c_1,\dots,c_k) = s_I(\sigma_1,\dots,\sigma_k)$ 
forms a basis of $\mathcal{S}^k$ as $I$ ranges over all partitions of $k$, 
by the injectivity of $\phi^*$, $s_I(c_1,\dots,c_k)$ also forms a basis of 
$H^{2k}(\mathrm{Gr}_n(\mathbb{C}^\infty);\mathbb{Z})$.

\subsection{A Product Formula}
Let $\omega$ be an $n$-dimensional complex vector bundle over
a compact complex manifold $B$. For each $k\geq0$ and partition 
$I=i_1,\dots,i_r$ of $k$, we define
\begin{equation*}
    s_I(c(\omega)) := s_I(c_1(\omega),\dots,c_k(\omega))\in
    H^{2k}(B;\mathbb{Z}).
\end{equation*}
sometimes we simply write it as $s_I(c)$. 

\begin{lemma}[Thom]
    \begin{equation*}
        s_I(c(\omega\oplus\omega')) =
        \sum_{J,K} s_J(c(\omega)) s_K(c(\omega')),
    \end{equation*}
    where the summation is taken over all partitions $J,K$
    with juxtapositon $JK$ equal to $I$.
\end{lemma}
\begin{remark}
    We can view this formula as a generalization of the Whitney sum formula.
\end{remark}
\begin{proof}
    Let $\phi:F(B)\to B$ be the flag manifold of $\omega\oplus\omega'$,
    with the splitting bundle
    \begin{equation*}
        \phi^*(\omega) = \eta_1\oplus\cdots\oplus\eta_n, \quad
        \phi^*(\omega') = \eta_{1}'\oplus\cdots\oplus\eta_{n'}'.
    \end{equation*}
    Then we have
    \begin{gather*}
        \phi^*c(\omega) = 
        (1+c_1(\eta_1))\cdots(1+c_1(\eta_n)) =: 
        (1+t_1)\cdots(1+t_n), \\
        \phi^*c(\omega') = 
        (1+c_1(\eta_1'))\cdots(1+c_1(\eta_{n'}')) =: 
        (1+t_{n+1})\cdots(1+t_{n+n'}).
    \end{gather*}
    Thus
    \begin{align*}
        \phi^*s_I(c(\omega\oplus\omega')) 
        &= \sum t_1^{i_1}\cdots t_r^{i_r} \\
        &= \sum_{
        \begin{subarray}{c}
            JK = I, \\ J = j_1,\dots,j_{r_1}, \\ K = k_1,\dots,k_{r_2}
        \end{subarray}
        } 
        \left(\sum t_1^{j_1}\cdots t_{r_1}^{j_{r_1}}\right)
        \left(\sum t_{n+1}^{k_1}\cdots t_{n+r_2}^{k_{r_2}}\right) \\
        &= \sum_{JK=I} s_J(c(\eta_1\oplus\cdots\oplus\eta_n)) s_K(c(\eta_1'\oplus\cdots\oplus\eta_n')) \\
        &= \sum_{JK=I} s_J(\phi^*c(\omega)) s_K(\phi^*c(\omega')) \\
        &= \phi^*\left(\sum_{JK=I} s_J(c(\omega)) s_K(c(\omega'))\right).
    \end{align*}
    Since $\phi^*$ is injective, we get the desired formula.
\end{proof}

Now for a compact complex manifold $K^n$. For each partition
$I=i_1,\dots,i_r$ of $n$, we denote
\begin{equation*}
    s_I[K^n] = s_I(c)[K^n] := \langle s_I(c(\tau)), \mu_{K^n}\rangle\in\mathbb{Z}.
\end{equation*}
This characteristic number is clearly a linear combination of
Chern numbers.

\begin{corollary}
    \begin{equation*}
        s_I[K^m\times L^n] =
        \sum_{JK=I} s_J[K^m] s_K[L^n].
    \end{equation*}
    where the summation is taken over all partitions $J$ of $m$ and $K$ of $n$
    with juxtapositon $JK$ equal to $I$.
\end{corollary}
\begin{proof}
    Since $\tau_{K\times L} = \tau_K \times \tau_L = \pi_1^*\tau_K \oplus \pi_2^*\tau_L$,
    where $\pi_1,\pi_2$ are the projections from $K\times L$ to
    $K$ and $L$ respectively, 
    \begin{align*}
        s_I[K\times L]
        &= \langle s_I(c(\tau_{K\times L})), \mu_{K\times L}\rangle \\
        &= \sum_{JK=I} \left\langle s_J(c(\pi_1^*\tau_K)) s_K(c(\pi_2^*\tau_L)),
        \mu_{K}\times \mu_{L} \right\rangle \\
        &= \sum_{JK=I} \langle s_J(c(\tau_K)), \mu_K \rangle
        \langle s_K(c(\tau_L)), \mu_L \rangle \\
        &= \sum_{JK=I} s_J[K] s_K[L].
    \end{align*}
\end{proof}

\begin{corollary}
    As a special case, $s_{m+n}[K^m\times L^n] = 0$, for $m,n>0$.
\end{corollary}

There are analogous results for Pontrjagin classes and numbers. 
Let $\xi$ be a real vector bundle over a compact smooth manifold $B$. 
For each $k\geq0$ and partition $I=i_1,\dots,i_r$ of $k$, we define
\begin{equation*}
    s_I(p(\xi)) := s_I(p_1(\xi),\dots,p_k(\xi))\in
    H^{4k}(B;\mathbb{Z}).
\end{equation*}
There are similar product formulas for Pontrjagin classes 
\begin{equation*}
    s_I(p(\xi\oplus\xi')) \equiv
    \sum_{J,K} s_J(p(\xi)) s_K(p(\xi')),
\end{equation*}
modulo elements of $2$-torsions. Since pairing with the fundamental class
annihilates all $2$-torsions, we have 
\begin{equation*}
    s_I[M\times N] = \sum_{JK=I} s_J[M] s_K[N],
\end{equation*}
Notes that these characteristic numbers
will be zero unless both the dimensions of $M$ and $N$ are divisible by $4$.

\subsection{Linear Independence of Chern Numbers and of Pontrjagin Numbers}
\begin{theorem}[Thom]
    Let $K^{1},\dots,K^{n}$ be compact complex manifolds
    with $s_k(c)[K^{k}]\neq0$. Then the $p(n)\times p(n)$ matrix
    \begin{equation*}
        \Big(c_{i_1}\cdots c_{i_r}[K^{j_1}\times\cdots\times K^{j_s}]\Big)
    \end{equation*}
    is nonsingular, where both $I=i_1,\dots,i_r$ and
    $J=j_1,\dots,j_s$ range over all partitions of $n$. 
\end{theorem} 
Similarly, we have the following theorem for Pontrjagin numbers.
\begin{theorem}[Thom]
    Let $M^{4},\dots,M^{4n}$ be compact oriented smooth manifolds
    with $s_k(p)[M^{4k}]\neq0$. Then the $p(n)\times p(n)$ matrix
    \begin{equation*}
        \Big(p_{i_1}\cdots p_{i_r}[M^{4j_1}\times\cdots\times M^{4j_s}]\Big)
    \end{equation*}
    is nonsingular, where both $I=i_1,\dots,i_r$ and
    $J=j_1,\dots,j_s$ range over all partitions of $n$.
\end{theorem}
\begin{proof}
    We only prove the first theorem, 
    the second one is completely analogous.

    Since the basis $c_{i_1}\cdots c_{i_r}$ isn't well behaved
    under the product operation, we change to the basis $s_I(c)$.
    By the product formula, 
    \begin{equation*}
        s_I[K^{j_1}\times\cdots\times K^{j_s}] \neq 0
    \end{equation*}
    only if $I=i_1,\dots,i_r$ is a refinement of $J=j_1,\dots,j_s$.
    Thus by suitable ordering of both the rows and columns, we can assume that
    the matrix is upper triangular. The diagonal elements are
    \begin{equation*}
        s_{i_1,\dots,i_r}[K^{i_1}\times\cdots\times K^{i_r}] =
        s_{i_1}[K^{i_1}]\cdots s_{i_r}[K^{i_r}] \neq 0.
    \end{equation*}
    Thus the matrix is nonsingular.
\end{proof}
% \section{Thom Space \& Transversality}
\label{sec:ThomSpace&Transversality}

\begin{definition}
    Let $\xi$ be a renk $k$ vector bundle over a base space $B$ with total space $E(\xi)$, 
    suppose $\xi$ is equipped with a Euclidean metric. Let 
    \begin{equation*}
        A = \{ v \in E(\xi) \mid \|v\| > 1 \}.
    \end{equation*}
    The \textbf{Thom space} of $\xi$ is defined to be the quotient space
    \begin{equation*}
        T(\xi) = E(\xi) / A.
    \end{equation*}
\end{definition}

\begin{lemma}\label{lem:CW-structure-of-Thom-space}
    If the base space $B$ is a CW-complex, then the Thom space $T(\xi)$ is 
    a $(k-1)$-connected CW-complex, having one cell of dimension $n+k$ for each
    cell of dimension $n$ in $B$ and a base point $t_0$.
\end{lemma}
\begin{remark}
    Notice that by the cellular approximation theorem, 
    a CW-complex with $k$-skeleton trivial is $(k-1)$-connected.
\end{remark}

\begin{lemma}\label{lem:Homology-of-Thom-space}
    If $\xi$ is an orientable rank $k$ vector bundle over $B$, then
    \begin{equation*}
        H_{k+i}(T(\xi),t_0;\mathbb{Z}) \cong H_i(B;\mathbb{Z}).
    \end{equation*}
\end{lemma}
\begin{proof}
    \begin{align*}
        H_{k+i}(T(\xi),t_0;\mZ) &= 
        H_{k+i}(E(\xi)/A,E_0/A;\mZ) \\
        &\cong H_{k+i}(E(\xi),E_0;\mZ) \quad \text{by excision thm}\\
        &\cong H_i(B;\mZ) \quad \text{by Thom isom}.
    \end{align*}
\end{proof}

\begin{theorem}\label{thm:C-isom-version-Hurewicz-thm}
    Let $X$ be a $(k-1)$-connected finite CW-complex, $k\geq2$. 
    Then the Hurewicz homomorphism
    \begin{equation*}
        \pi_r(X)\to H_r(X;\mathbb{Z})
    \end{equation*}
    is a $\mathcal{C}$-isomorphism for $r<2k-1$. 
\end{theorem}

\begin{corollary}
    If $T$ is the Thom space of an \textcolor{red}{orientable} rank $k$ vector bundle over a
    \textcolor{red}{finite} CW-complex $B$, $k\geq2$, then there is a $\mathcal{C}$-isomorphism
    \begin{equation*}
        \pi_{k+i}(T) \to H_i(B;\mathbb{Z})
    \end{equation*}
    for $i<k-1$.
\end{corollary}
\begin{proof}
    \begin{equation*}
        \pi_{k+i}(T) \overset{\mathcal{C}-\text{isom}}{\cong} H_{k+i}(T,t_0;\mZ) \cong H_i(B;\mZ), \quad i<k-1.
    \end{equation*}
    The first isomorphism holds since by Lemma \ref{lem:CW-structure-of-Thom-space} 
    $T$ is $(k-1)$-connected finite CW-complex, then we can apply Theorem \ref{thm:C-isom-version-Hurewicz-thm};
    the second isomorphism holds since $\xi$ is orientable, we can apply Lemma \ref{lem:Homology-of-Thom-space}.
\end{proof}

\begin{theorem}
    Every continuous map $f\colon S^{m} \to B$ is homotopic to a map 
    $g$ which is smooth throughout $g^{-1}(T-t_0)$ and transverse to the zero section $B$.
    The oriented cobordism class of the submanifold $g^{-1}(B)$ depends only on the homotopy 
    class of $g$. Hence there is a well-defined map
    \begin{align*}
        \pi_{m}(T,t_0) &\to \Omega^{\mathrm{SO}}_{m-k} \\
        [g] &\mapsto [g^{-1}(B)].
    \end{align*}
\end{theorem}

\subsection{The Main Theorem}
Let $\tilde{\gamma}^k$ be the universal oriented rank 
$k$ vector bundle over $\tilde{\mrm{Gr}}_k(\mR^\infty)$, 
we have the following important theorem.
\begin{theorem}[Thom's Theorem]
    For $k>n+1$, the map
    \begin{align*}
        \pi_{n+k}(T(\tilde{\gamma}^k),t_0) &\to \Omega^{\mrm{SO}}_{n} \\
        [g] &\mapsto [g^{-1}(\tilde{\mrm{Gr}}_k(\mR^\infty))]
    \end{align*}
    is an isomorphism. Similarly, for the unoriented case,
    the map
    \begin{align*}
        \pi_{n+k}(T(\tilde{\gamma}^k),t_0) &\to \Omega_{n} \\
        [g] &\mapsto [g^{-1}(\tilde{\mrm{Gr}}_k(\mR^\infty))]
    \end{align*}
    is also an isomorphism.
\end{theorem}
\begin{remark}
    这个同构是否由这个映射给出存疑.
\end{remark}
We  will not prove this theorem, instead, we will proof part of it.
Let $\tilde{\gamma}^k_p$ be the universal bundle over $\tilde{\mrm{Gr}}_k(\mR^{k+p})$.
\begin{lemma}
    If $k\geq n$ and $p\geq n$, then the homomorphism
    \begin{align*}
        \pi_{n+k}(T(\tilde{\gamma}^k_p),t_0) &\to \Omega^{\mrm{SO}}_{n} \\
        [g] &\mapsto [g^{-1}(\tilde{\mrm{Gr}}_k(\mR^{k+p}))]
    \end{align*}
    is surjective.
\end{lemma}
\section{Multiplicative sequences}
\label{sec:multiplicative_sequences}

We start with polynomial ring $\mQ[\beta_1,\beta_2,\dots]$ in countably many variables $\beta_i$ over the rational number field $\mQ$. 
Let $A=\mQ[a_1,a_2,\dots]=\mathrm{Sym}[\beta_1,\beta_2,\dots]$ be the subring of symmetric polynomials, 
where $a_i$ is the $i$-th elementary symmetric polynomial in $\beta_1,\beta_2,\dots$. $A$ has a grading structure by setting $\deg a_i=i$.
\begin{equation*}
    A=\bigoplus_{n=0}^{\infty} A^n, \quad A^n=\big\{\text{homogeneous polynomials of degree } n\big\}.
\end{equation*}
We set
\begin{equation*}
    a = 1 + a_1t + a_2t^2 + \cdots = \prod_{i=1}^{\infty} (1+\beta_it) \in A[[t]],
\end{equation*}
For any two 
\begin{gather*}
    A = \mQ[a_1,a_2,\dots]\subset \mQ[\beta_1,\beta_2,\dots], \quad B = \mQ[b_1,b_2,\dots]\subset \mQ[\beta_1',\beta_2',\dots]
\end{gather*}
and
\begin{gather*}
    a = 1 + a_1t + a_2t^2 + \cdots, \quad b = 1 + b_1t + b_2t^2 + \cdots
\end{gather*}
their product is 
\begin{align*}
    ab =& (1 + a_1t + a_2t^2 + \cdots)(1 + b_1t + b_2t^2 + \cdots) \\
    =& 1 + (a_1 + b_1)t + (a_2 + a_1b_1 + b_2)t^2 + \cdots
\end{align*}

\begin{remark}
    If we take $t=1$ and $a_i = c_i(\omega)$, where $\omega$ is a complex $n$-plane bundle, then 
    \begin{equation*}
        a = 1 + c_1(\omega) + c_2(\omega) + \cdots + c_n(\omega) = c(\omega)
    \end{equation*}
\end{remark}

\subsection{Multiplicative Sequences}
Let $\{K_n\}$ be a sequence of polynomials:
\begin{equation*}
    K_0 = 1, \quad K_1 = K_1(x_1), \quad K_2 = K_2(x_1,x_2), \quad K_3 = K_3(x_1,x_2,x_3), \quad \cdots,
\end{equation*} 
satisfying that each 
\begin{equation*}
    K_n(a_1,\dots,a_n)\in A^n
\end{equation*}
is a homogeneous polynomial of degree $n$. 

For any 
\begin{equation*}
    a = 1 + a_1t + a_2t^2 + \cdots, 
\end{equation*}
we define
\begin{equation*}
    K(a) = 1 + K_1(a_1)t + K_2(a_1,a_2)t^2 + K_3(a_1,a_2,a_3)t^3 + \cdots.
\end{equation*}
We say that $\{K_n\}$ is a \textbf{multiplicative sequence} (or m-sequence) if for any two such $a$ and $b$ as above, we have 
\begin{equation*}
    K(ab) = K(a)K(b) 
\end{equation*}
which forces that
\begin{align*}
    &K_1(a_1+b_1) = K_1(a_1) + K_1(b_1), \\
    &K_2(a_1+b_1,a_2+a_1b_1+b_2) = K_2(a_1,a_2) + K_1(a_1)K_1(b_1) + K_2(b_1,b_2), \\
    &\cdots\cdots
\end{align*}

\begin{example}\label{ex:m-sequence1}
    Set $K_n(x_1,\dots,x_n) = \lambda^n x_n$, where $\lambda$ is a constant. Then we have
    \begin{equation*}
        K(1+a_1t+a_2t^2+\cdots) = 1 + \lambda a_1 t + \lambda^2 a_2 t^2 + \cdots 
    \end{equation*}
    Especially, if $\lambda = 1$, then 
    \begin{equation*}
        K(a) = a.
    \end{equation*}
    If $\lambda = -1$, then 
    \begin{equation*}
        K(a) = 1 - a_1t + a_2t^2 - a_3t^3 + \cdots. 
    \end{equation*}
    If in addition we take $t=1$, $a_i = c_i(\omega)$, then we get
    \begin{equation*}
        K(c(\omega)) = 1 - c_1(\omega) + c_2(\omega) - \cdots + (-1)^n c_n(\omega) = c(\bar{\omega}).
    \end{equation*}
\end{example}
\begin{example}\label{ex:m-sequence2}
    The identity $K(a) = a^{-1}$ defines a multiplicative sequence, where
    \begin{align*}
        &K_0 = 1, \\
        &K_1(x_1) = -x_1, \\
        &K_2(x_1,x_2) = x_1^2 - x_2, \\
        &K_3(x_1,x_2,x_3) = -x_1^3 + 2x_1x_2 - x_3, \\
        &K_4(x_1,x_2,x_3,x_4) = x_1^4 - 3x_1^2x_2 + 2x_1x_3 + x_2^2 - x_4, \\
        &\cdots\cdots
    \end{align*}
    If we take $t=1$, $a_i = c_i(\omega)$ and $\eta\oplus\omega = \varepsilon^{m+n}$, then we get
    \begin{equation*}
        K(c(\omega)) = c(\omega)^{-1} = c(\eta).
    \end{equation*}
\end{example}
\begin{example}\label{ex:m-sequence3}
    The polynomials 
    \begin{align*}
        &K_{2n+1} = 0, \\
        &K_{2n} = x_n^2 - 2x_{n-1}x_{n+1} + \cdots + (-1)^{n-1}2x_1 x_{2n-1} + (-1)^n 2n x_{2n}
    \end{align*}
    form an m-sequence. If we take $t=1$, $a_i = c_i(\omega)$, then we get the total Pontryagin class
    \begin{equation*}
        K(c(\omega)) = 1 + p_1(\omega_{\mR}) + p_2(\omega_{\mR}) + \cdots = p(\omega_{\mR}).
    \end{equation*}
\end{example}

\subsection{Characteristic Power Series}
Given a multiplicative sequence $\{K_n\}$, we can define its \textbf{characteristic power series} $f(t)$ by 
\begin{equation*}
    K(1 + t) = f(t) = 1 + \lambda_1 t + \lambda_2 t^2 + \lambda_3 t^3 + \cdots,
\end{equation*}
where $\lambda_i = K_i(1,0,\dots,0)$. The following lemma shows that there is a one-to-one correspondence 
between multiplicative sequences and power series.
\begin{lemma}[Hirzebruch]
    Let $f(t) = 1 + \lambda_1 t + \lambda_2 t^2 + \lambda_3 t^3 + \cdots$ be a power series in $\mQ[[t]]$ with constant term 1. 
    Then there \textbf{exists} a \textbf{unique} multiplicative sequence $\{K_n\}$ whose characteristic power series is $f(t)$.
\end{lemma}
\begin{proof}
    Uniqueness: Suppose that $\{K_n\}$ is a multiplicative sequence with characteristic power series $f(t)$. 
    Then 
    \begin{equation*}
        K(1 + a_1 t + a_2 t^2 + \cdots) = K\left(\prod_{i=1}^{\infty} (1 + \beta_i t)\right) 
        = \prod_{i=1}^{\infty} K(1 + \beta_i t) = \prod_{i=1}^{\infty} f(\beta_i t).
    \end{equation*}
    take $\beta_{n+1} = \beta_{n+2} = \cdots = 0$, we have
    \begin{equation*}
        K(1 + a_1 t + a_2 t^2 + \cdots + a_n t^n) = \prod_{i=1}^{n} f(\beta_i t).
    \end{equation*}
    by comparing the coefficient of $t^n$ on both sides, 
    \begin{equation*}
        K_n(a_1,a_2,\dots,a_n) = \text{the coefficient of } t^n \text{ in } \prod_{i=1}^{n} f(\beta_i t).
    \end{equation*}
    Since the right-hand side is a symmetric polynomial in $\beta_1,\dots,\beta_n$, 
    it can be expressed as a polynomial in $a_1,\dots,a_n$. 
    And $a_1,\dots,a_n$ are algebraically independent, so $K_n$ is uniquely determined.

    Existence: Define $K_n$ by the above formula, then it is easy to check that $\{K_n\}$ is a multiplicative sequence
    whose characteristic power series is $f(t)$.
\end{proof}
\begin{remark}
    Given a power series 
    \begin{equation*}
        f(t) = 1 + \lambda_1 t + \lambda_2 t^2 + \lambda_3 t^3 + \cdots,
    \end{equation*}
    the corresponding multiplicative sequence $\{K_n\}$ can be computed by the following: 
    Denote $a_i^{(n)} = a_i(\beta_1,\dots,\beta_n,0,0,\dots)$, then
    \begin{align*}
        K_1(a_1^{(1)}) &= \text{the coefficient of } t \text{ in } f(\beta_1 t) \\
        & = \lambda_1 \beta_1 = \lambda_1 a_1^{(1)} \\
        K_2(a_1^{(2)},a_2^{(2)}) &= \text{the coefficient of } t^2 \text{ in } f(\beta_1 t) f(\beta_2 t) \\
        & = \lambda_1^2 \beta_1\beta_2 + \lambda_2 (\beta_1^2 + \beta_2^2), \\
        & = \lambda_1^2 a_2^{(2)} + \lambda_2 \big((a_1^{(2)})^2 - 2a_2^{(2)}\big) \\
        K_3(a_1^{(3)},a_2^{(3)},a_3^{(3)}) &= \text{the coefficient of } t^3 \text{ in } f(\beta_1 t) f(\beta_2 t) f(\beta_3 t) \\
        & = \lambda_1^3 \beta_1\beta_2\beta_3 + \lambda_1\lambda_2 (\beta_1\beta_2^2 + \beta_2\beta_3^2 + \beta_3\beta_1^2) + \lambda_3 (\beta_1^3 + \beta_2^3 + \beta_3^3) \\
        & = \lambda_1^3 a_3^{(3)} + \lambda_1\lambda_2 \big(a_1^{(3)} a_2^{(3)} - 3 a_3^{(3)}\big) + \lambda_3 \big((a_1^{(3)})^3 - 3 a_1^{(3)} a_2^{(3)} + 3 a_3^{(3)}\big) \\
        & \cdots\cdots
    \end{align*}
\end{remark}
\begin{example}
    The characteristic power series of the multiplicative sequence in Example \ref{ex:m-sequence1} is $f(t) = 1 + \lambda t$.
\end{example}
\begin{example}
    The characteristic power series of the multiplicative sequence in Example \ref{ex:m-sequence2} is 
    \begin{equation*}
        f(t) = \frac{1}{1+t} = 1 - t + t^2 - t^3 + t^4 - \cdots.
    \end{equation*}
    We can check the first few polynomials:
    \begin{align*}
        &K_0 = 1, \\
        &K_1(x_1) = \lambda_1 x_1 = -x_1, \\
        &K_2(x_1,x_2) = \lambda_1^2 x_2 + \lambda_2 (x_1^2 - 2x_2) = x_1^2 - x_2, \\
        &K_3(x_1,x_2,x_3) = \lambda_1^3 x_3 + \lambda_1\lambda_2 (x_1 x_2 - 3 x_3) + \lambda_3 (x_1^3 - 3 x_1 x_2 + 3 x_3) = -x_1^3 + 2 x_1 x_2 - x_3, \\[4pt]
        &\cdots\cdots
    \end{align*}
\end{example}
\begin{example}
    The characteristic power series of the multiplicative sequence in Example \ref{ex:m-sequence3} is $f(t) = 1 + t^2$.
    We can check the first few polynomials:
    \begin{align*}
        &K_0 = 1, \\
        &K_1(x_1) = \lambda_1 x_1 = 0, \\
        &K_2(x_1,x_2) = \lambda_2 x_2 = x_2, \\
        &K_3(x_1,x_2,x_3) = \lambda_1^3 x_3 + \lambda_1\lambda_2 (x_1 x_2 - 3 x_3) + \lambda_3 (x_1^3 - 3 x_1 x_2 + 3 x_3) = 0, \\
        &K_4(x_1,x_2,x_3,x_4) = \lambda_2^2 (x_2^2 - 2 x_1 x_3 + 2 x_4) = x_2^2 - 2 x_1 x_3 + 2 x_4, \\[4pt]
        &\cdots\cdots
    \end{align*}
\end{example}

\subsection{K-genus}
\begin{definition}
    Given a multiplicative sequence $\{K_n\}$. For any smooth, closed, oriented manifold $M^m$, 
    we can define its \textbf{K-genus} by 
    \begin{align*}
        &K[M] = 0, \quad \text{if $4 \nmid m$}, \\
        &K[M] = \langle K_{n}(p_1(\tau_M), \dots, p_n(\tau_M)), \mu_M \rangle, \quad \text{if } m=4n,
    \end{align*}
    Thus $K[M]$ is a certain rational linear combination of Pontryagin numbers of $M$.
\end{definition}
\begin{remark}
    A K-genus can be written briefly as 
    \begin{equation*}
        K[M] = \langle K(p(\tau_M)), \mu_M \rangle,
    \end{equation*}
\end{remark}
\begin{lemma}
    A K-genus gives rise to a ring homomorphism 
    \begin{equation*}
        \Omega_*^{SO} \to \mQ,\qquad [M] \mapsto K[M]
    \end{equation*}
    from the oriented cobordism ring to the rational number field. 

    Equivalently, this correspondence gives an algebra homomorphism from $\Omega_*^{SO} \otimes \mQ$ to $\mQ$. 
    Thus defines an element in the dual space $\mathrm{Hom}_{\mQ}(\Omega_*^{SO} \otimes \mQ, \mQ)$.
\end{lemma}
\begin{proof}
    The correspondence is well-defined since Pontryagin numbers are cobordism invariants.
    It is clear that $K$ is additive under disjoint union.
    
    It suffices to show that for any two $M$ and $M'$, we have 
    \begin{equation*}
        K[M \times M'] = K[M] K[M'].
    \end{equation*}
    If either $\dim M$ or $\dim M'$ is not divisible by 4, then both sides are zero. 
    Now suppose that $\dim M = 4m$ and $\dim M' = 4n$. Then 
    \begin{align*}
        K[M \times M'] =& \langle K(p(\tau_{M \times M'})), \mu_{M \times M'} \rangle \\
        =& \langle K(p(\tau_M) \times p(\tau_{M'})), \mu_M \times \mu_{M'} \rangle \\
        =& \langle K(p(\tau_M))\times K(p(\tau_{M'})), \mu_M \times \mu_{M'} \rangle \\
        =& \langle K(p(\tau_M)), \mu_M \rangle \cdot \langle K(p(\tau_{M'})), \mu_{M'} \rangle \\
        =& K[M] K[M'].
    \end{align*}
\end{proof}

\begin{remark}
    The multiplicative sequence and K-genus give a systematic machinery to produce various invariants of manifolds.
\end{remark}

\subsection{Signature}
\begin{definition}
    The \textbf{signature} $\sigma(M)$ of a smooth, closed, oriented $4k$-manifold $M$ is defined to be the signature of the symmetric bilinear form
    \begin{equation*}
        Q_M : H^{2k}(M;\mQ) \times H^{2k}(M;\mQ) \to \mQ, \quad Q_M(a,b) = \langle a \cup b, \mu_M \rangle.
    \end{equation*}
    For $M$ with $4\nmid \dim M$, we set $\sigma(M) = 0$.
\end{definition}
To be precise, if we choose a basis $\{a_1,\dots,a_r\}$ of $H^{2k}(M;\mQ)$, so that the matrix of $Q_M$
\begin{equation*}
    Q_{ij} = \langle a_i \cup a_j, \mu_M \rangle,
\end{equation*}
is diagonal. Then the signature $\sigma(M)$ is defined to be the number of positive entries minus the number of negative entries on the diagonal. 
Notes that $Q_M$ is non-degenerate by Poincar\'e duality, so the matrix $(Q_{ij})$ is invertible. 

\begin{lemma}[Thom]
    The signature $\sigma(M)$ gives rise to a ring homomorphism
    \begin{equation*}
        \Omega_*^{SO} \to \mZ, \quad [M] \mapsto \sigma(M).
    \end{equation*}
    Equivalently, this is to say the signature function satisfies:
    \begin{enumerate}
        \item $\sigma(M \sqcup M') = \sigma(M) + \sigma(M')$,
        \item $\sigma(M \times M') = \sigma(M) \sigma(M')$.
        \item If $M = \p W$ for some oriented manifold $W$, then $\sigma(M) = 0$.
    \end{enumerate}
\end{lemma}
\begin{proof}
    The first property follows by 
    \begin{equation*}
        H^{2k}(M \sqcup M';\mQ) \cong H^{2k}(M;\mQ) \oplus H^{2k}(M';\mQ).
    \end{equation*}
    The second property follows by the K\"unneth formula
    \begin{equation*}
        H^*(M \times M';\mQ) \cong H^*(M;\mQ) \otimes H^*(M';\mQ).
    \end{equation*}
    For the third property, we need to show that if $M = \p W$, then $\sigma(M) = 0$.
    By the long exact sequence of the pair $(W,M)$ and Poincar\'e-Lefschetz duality, we have the following commutative diagram:
    \begin{equation*}
        \begin{tikzcd}
            H^{2k}(W,M;\mQ) \arrow[r, "i^*"] \arrow[d, "\cong"] & H^{2k}(W;\mQ) \arrow[r, "j^*"] \arrow[d, "\cong"] & H^{2k}(M;\mQ) \arrow[d, "\cong"] \\
            H_{4k-2k}(W;\mQ) \arrow[r, "j_*"] & H_{4k-2k}(W,M;\mQ) \arrow[r, "i_*"] & H_{4k-2k-1}(M;\mQ)
        \end{tikzcd}
    \end{equation*}
    Thus we have an exact sequence
    \begin{equation*}
        H^{2k}(W,M;\mQ) \xrightarrow{i^*} H^{2k}(W;\mQ) \xrightarrow{j^*} H^{2k}(M;\mQ) \xrightarrow{\delta} H^{2k+1}(W,M;\mQ).
    \end{equation*}
    Note that for any $a,b\in H^{2k}(M;\mQ)$, if either $a$ or $b$ is in $\mathrm{Im}\, j^*$, then 
    \begin{equation*}
        Q_M(a,b) = \langle a\cup b, \mu_M \rangle = 0,
    \end{equation*}
    since $j_* (\mu_M) = 0$. Thus $Q_M$ descends to a well-defined bilinear form on 
    \begin{equation*}
        V = H^{2k}(M;\mQ)/\mathrm{Im}\, j^*.
    \end{equation*}
    Moreover, by the exactness of the above sequence, we have an isomorphism
    \begin{equation*}
        V \cong \mathrm{Im}\, \delta \subset H^{2k+1}(W,M;\mQ).
    \end{equation*}
    Since $\dim H^{2k+1}(W,M;\mQ)$ is even, $\dim V$ is also even.
    On the other hand, the bilinear form $Q_M$ on $V$ is skew-symmetric, i.e. $Q_M([a],[b]) = - Q_M([b],[a])$ for any $[a],[b]\in V$.
    Thus the signature of $Q_M$ on $V$ is zero. Since $Q_M$ vanishes on $\mathrm{Im}\, j^*$, we have
    \begin{equation*}
        \sigma(M) = \sigma(Q_M) = \sigma(Q_M|_V) = 0.
    \end{equation*}
\end{proof}
\begin{remark}
    This lemma shows that the signature determine an element in the dual space $\mathrm{Hom}_{\mQ}(\Omega_*^{SO}\otimes \mQ, \mQ)$ 
    which takes values in $\mZ$ for each oriented manifold. 
\end{remark}
Since 
\begin{equation*}
    \Omega_*^{SO} \otimes \mQ \cong \mQ[\mC P^2, \mC P^4, \dots]
\end{equation*}
is a polynomial ring generated by the complex projective spaces, 
and the dual space 
\begin{equation*}
    \mathrm{Hom}_{\mQ}(\Omega_*^{SO}\otimes \mQ, \mQ)\cong \mrm{span}_{\mQ}\{p_I\}_I
\end{equation*}
is spanned by the Pontryagin numbers. So the signature function can be written as a linear combination of Pontryagin numbers. 

\begin{theorem}[Signature Theorem, Hirzebruch]
    The signature $\sigma(M)$ of a smooth, closed, oriented $4k$-manifold $M$ is equal to the $L$-genus of $M$, i.e.
    \begin{equation*}
        \sigma(M) = \langle L_k(p_1(\tau_M), \dots, p_k(\tau_M)), \mu_M \rangle,
    \end{equation*}
    where $L_k$ is the multiplicative sequence with characteristic power series 
    \begin{equation*}
        f(t) = \frac{\sqrt{t}}{\tanh \sqrt{t}} = 1 + \frac{t}{3} - \frac{t^2}{45} + \frac{2 t^3}{945} - \cdots + 
        (-1)^{k-1} \frac{2^{2k} B_{2k}}{(2k)!} t^k + \cdots,
    \end{equation*}
\end{theorem}
Here $B_{2k}$ are the Bernoulli numbers defined by the generating function
\begin{equation*}
    \frac{t}{e^t - 1} = 1 - \frac{t}{2} + \sum_{k=1}^{\infty} B_{2k} \frac{t^{2k}}{(2k)!}.
\end{equation*}
The first few Bernoulli numbers are
\begin{equation*}
    B_0 = 1, \quad B_1 = \frac{1}{6}, \quad B_2 = \frac{1}{30}, \quad B_3 = \frac{1}{42}, \quad B_4 = \frac{1}{30}, \quad B_5 = \frac{5}{66}, \quad \cdots
\end{equation*}
The first few $L$-polynomials are
\begin{align*}
    &L_0 = 1, \\
    &L_1 = \frac{1}{3} p_1, \\
    &L_2 = \frac{1}{45} (7 p_2 - p_1^2), \\
    &L_3 = \frac{1}{945} (62 p_3 - 13 p_1 p_2 + 2 p_1^3), \\
    &L_4 = \frac{1}{14175} (381 p_4 - 71 p_1 p_3 - 19 p_2^2 + 22 p_1^2 p_2 - 3 p_1^4), \\[6pt]
    &\cdots\cdots
\end{align*} 
Here we compute the first three $L$-polynomials as follows:
\begin{align*}
    &L_1(p_1) = \lambda_1 p_1 = \frac{1}{3} p_1, \\
    &L_2(p_1,p_2) = \lambda_1^2 p_2 + \lambda_2 (p_1^2 - 2 p_2) = \frac{1}{9} p_2 - \frac{1}{45} (p_1^2 - 2 p_2) = \frac{7}{45} p_2 - \frac{1}{45} p_1^2, \\
    &L_3(p_1,p_2,p_3) = \lambda_1^3 p_3 + \lambda_1\lambda_2 (p_1 p_2 - 3 p_3) + \lambda_3 (p_1^3 - 3 p_1 p_2 + 3 p_3) \\
    & = \frac{1}{27} p_3 - \frac{1}{135} (p_1 p_2 - 3 p_3) + \frac{2}{945} (p_1^3 - 3 p_1 p_2 + 3 p_3) \\
    & = \frac{62}{945} p_3 - \frac{13}{945} p_1 p_2 + \frac{2}{945} p_1^3.
\end{align*}
\begin{proof}[Proof of the Signature Theorem]
    Since the correspondences
    \begin{gather*}
        \Omega_*^{SO} \otimes \mQ \to \mQ, \quad [M] \mapsto \sigma(M) \\
        \Omega_*^{SO} \otimes \mQ \to \mQ, \quad [M] \mapsto L[M]
    \end{gather*}
    are both ring homomorphisms, it suffices to check that they agree on the polynomial generators 
    $\mC P^{2}, \mC P^{4}, \dots$ of $\Omega_*^{SO} \otimes \mQ$.

    \noindent \textbf{Step 1:} Compute $\sigma(\mC P^{2k})$.
    Since $H^k(\mC P^{2k};\mQ)$ is generated by $a^k$ with 
    \begin{equation*}
        \langle a^k\smile, \mu_{\mC P^{2k}} \rangle = 1,
    \end{equation*} 
    Thus $\sigma(\mC P^{2k}) = 1$. 

    \noindent \textbf{Step 2:} Compute $L[\mC P^{2k}]$.
    The total Pontryagin class of $\mC P^{2k}$ is equal to
    \begin{equation*}
        p = p(\tau_{\mC P^{2k}}) = (1 + a^2)^{2k+1},
    \end{equation*}
    Since the characteristic power series of $L$ is
    \begin{equation*}
        f(t) = \frac{\sqrt{t}}{\tanh \sqrt{t}} 
    \end{equation*}
    We have
    \begin{align*}
        L(p) =& \big[f(a^2)\big]^{2k+1} = \left(\frac{a}{\tanh a}\right)^{2k+1} \\
        =& \left(1 - \frac{a^2}{3} + \frac{2 a^4}{15} - \frac{17 a^6}{315} + \cdots\right)^{2k+1}.
    \end{align*}
    Thus 
    \begin{equation*}
        L[\mC P^{2k}] = \langle L(p), \mu_{\mC P^{2k}} \rangle = \text{the coefficient of } a^{2k} \text{ in } \left(\frac{a}{\tanh a}\right)^{2k+1}.
    \end{equation*}
    In order to compute this coefficient, we use the residue theorem in complex analysis: let 
    \begin{equation*}
        g(z) = \left(\frac{z}{\tanh z}\right)^{2k+1}
    \end{equation*}
    be a complex function. Then the coefficient of $z^{2k}$ in $g(z)$ is equal to
    \begin{align*}
        \frac{1}{2\pi i}\int_{|z|=\varepsilon} \frac{g(z)}{z^{2k+1}} \md z &= \frac{1}{2\pi i}\int_{|z|=\varepsilon} \frac{\md z}{(\tanh z)^{2k+1}}  \\
        &= \frac{1}{2\pi i}\int_{|u|=\varepsilon'} \frac{\md u}{(1-u^2)u^{2k+1}} \quad (u = \tanh z) \\ 
        &= \frac{1}{2\pi i}\int_{|u|=\varepsilon'} \frac{1+u^2+u^4+\cdots}{u^{2k+1}} \md u \\
        &= 1,
    \end{align*}
    where $\varepsilon$ and $\varepsilon'$ are sufficiently small. Thus we have $L[\mC P^{2k}] = 1$.
\end{proof}
A more direct proof of the signature theorem is given by Atiyah and Singer in their famous index theorem for elliptic operators.

\begin{corollary}
    The L-genus $L[M]$ of any smooth, closed, oriented $4k$-manifold $M$ is an integer.
\end{corollary}
\begin{proof}
    Since $L[M] = \sigma(M) \in \mZ$.
\end{proof}
It follows that the Pontryagin numbers of any smooth, closed, oriented $4k$-manifold satisfy certain integrality condition. 
For example, 
\begin{align*}
    &L_1 = \frac{1}{3} p_1\Rightarrow p_1[\tau_M] \equiv 0 \mod 3, \\
    &L_2 = \frac{1}{45} (7 p_2 - p_1^2) \Rightarrow 7 p_2[\tau_M] - p_1^2[\tau_M] \equiv 0 \mod 45.
\end{align*}

\begin{corollary}
    The L-genus $L[M]$ is an oriented homotopy invariant of $M$.
\end{corollary}
\begin{proof}
    Since $\sigma(M)$ is invariant under orientation-preserving homotopy equivalences. 
\end{proof}

\subsection{Exercise}
\begin{problemstar}[19-A]
Let ${T_n}$ be the multiplicative sequence of polynomials belonging to the power series
$f(t) = \frac{t}{1-\me^{-t}}$. Then the \textbf{Todd genus} $T[M]$ of a complex $n$-dimensional
manifold is defined to be the characteristic number $\langle T_n(c_1(\tau_M), \dots, c_n(\tau_M)), \mu_M \rangle$. 
Prove that $T[\mC P^n] = +1$ and prove that ${T_n}$ is the only multiplicative sequence with this property.
\end{problemstar}




% \bibliography{References}
\end{document}

